\chapter{Functions and Variations}\label{ch:g11:func}

This chapter organizes what one needs to know about \omterm{def:g11:func:function}{functions} before
calculus: the small family of \emph{\omterm{def:g11:func:reference}{reference functions}} whose \omterm{def:g10:functions:graph}{graphs} must
be known by heart, the precise meaning of \emph{\omterm{def:g11:func:monotone}{increasing}} and
\emph{\omterm{def:g10:functions:variations}{decreasing}}, and how variations behave when \omterm{def:g11:func:function}{functions} are combined.

\section{Functions and curves}

\begin{definition}[Function, domain, curve]\label{def:g11:func:function}
A \emph{function}\index{function} $f$ assigns to each number $x$ of a set
$D \subseteq \R$ (its \emph{domain}\index{domain}) exactly one number
$f(x)$. Its \emph{curve} (or \omterm{def:g10:functions:graph}{graph}) is the set of points
$\bigl(x, f(x)\bigr)$ for $x \in D$.
\end{definition}

\begin{example}
The formula $f(x) = \frac{\sqrt{x+2}}{x-1}$ requires $x + 2 \geq 0$ (for
the square root) and $x \neq 1$ (nonzero denominator): the \omterm{def:g11:func:function}{domain} is
$\intco{-2}{1} \cup \intoo{1}{+\infty}$.
\end{example}

\section{The reference functions}

The following five \omterm{def:g11:func:function}{functions} and their curves must become reflexes.

\begin{definition}[Reference functions]\label{def:g11:func:reference}
\index{reference functions}
\[
x \mapsto x^2 \text{ (square)}, \quad
x \mapsto x^3 \text{ (cube)}, \quad
x \mapsto \sqrt{x}, \quad
x \mapsto \frac{1}{x}, \quad
x \mapsto \abs{x}.
\]
Their \omterm{def:g11:func:function}{domains} are $\R$, $\R$, $\intco{0}{+\infty}$, $\R \setminus \{0\}$
and $\R$ respectively. The curve of $x \mapsto \frac1x$ is a
\emph{hyperbola}\index{hyperbola}.
\end{definition}

\begin{omfigure}
\begin{tikzpicture}
\begin{axis}[omaxis,
  width=7cm, height=5.6cm,
  xmin=-2.3, xmax=2.3, ymin=-3.4, ymax=3.6,
  xtick={-1,1}, ytick={-1,1}]
  \addplot[omDef, thick, domain=-1.85:1.85] {x^2};
  \addplot[omThm, thick, domain=-1.5:1.5] {x^3};
  \node[omDef, font=\small, anchor=east] at (axis cs:-1.5,1.9) {$x^2$};
  \node[omThm, font=\small, anchor=east] at (axis cs:-1.32,-1.73) {$x^3$};
\end{axis}
\end{tikzpicture}%
\begin{tikzpicture}
\begin{axis}[omaxis,
  width=7cm, height=5.6cm,
  xmin=-3.4, xmax=4.4, ymin=-3.4, ymax=3.6,
  xtick={1}, ytick={1}]
  \addplot[omProp, thick, domain=-3.2:-0.31, samples=60] {1/x};
  \addplot[omProp, thick, domain=0.31:4.2, samples=60] {1/x};
  \addplot[omDef, thick, domain=0:4.2] {sqrt(x)};
  \addplot[omThm, thick, domain=-3.2:3.2] {abs(x)};
  \node[omDef, font=\small, anchor=north] at (axis cs:3.6,1.55)
    {$\sqrt{x}$};
  \node[omThm, font=\small, anchor=east] at (axis cs:-2.4,2.8)
    {$\abs{x}$};
  \node[omProp, font=\small, anchor=west] at (axis cs:0.75,2.4)
    {$\tfrac1x$};
\end{axis}
\end{tikzpicture}

{\small The five reference curves: the \omterm{def:g11:quad:parabola}{parabola} $x^2$ and the cube $x^3$
(left); $\sqrt x$, $\abs x$ and the \omterm{def:g11:func:reference}{hyperbola} $\frac1x$ (right).}
\end{omfigure}

\section{Variations on an interval}

\begin{definition}[Increasing, decreasing]\label{def:g11:func:monotone}
Let $f$ be defined on an \omterm{def:g10:numbers:interval}{interval} $I$. The \omterm{def:g11:func:function}{function} $f$ is
\emph{increasing}\index{increasing function} on $I$ if it preserves order:
\[
\text{for all } u, v \in I: \quad u < v \implies f(u) \leq f(v),
\]
and \emph{\omterm{def:g10:functions:variations}{decreasing}} if it reverses order ($u < v \implies f(u) \geq
f(v)$). With strict inequalities on the right, $f$ is \emph{strictly}
increasing (resp.\ \omterm{def:g10:functions:variations}{decreasing}). A \omterm{def:g11:func:function}{function} that is increasing or
\omterm{def:g10:functions:variations}{decreasing} on $I$ is \emph{monotonic} on $I$.
\end{definition}

\begin{method}[Proving a variation from the definition]\label{met:g11:func:variation}
Take two arbitrary numbers $u < v$ in the \omterm{def:g10:numbers:interval}{interval}, and determine the sign
of the difference $f(v) - f(u)$, usually by \emph{\omterm{def:g10:algebra:expand}{factoring}} it. If the
sign is constant on the \omterm{def:g10:numbers:interval}{interval}, conclude.
\end{method}

\begin{proposition}[Variations of the reference functions]\label{prop:g11:func:refvariations}
\begin{enumerate}
  \item $x^2$ is strictly \omterm{def:g10:functions:variations}{decreasing} on $\intoc{-\infty}{0}$ and strictly
        \omterm{def:g11:func:monotone}{increasing} on $\intco{0}{+\infty}$.
  \item $x^3$ is strictly \omterm{def:g11:func:monotone}{increasing} on $\R$.
  \item $\sqrt x$ is strictly \omterm{def:g11:func:monotone}{increasing} on $\intco{0}{+\infty}$.
  \item $\frac1x$ is strictly \omterm{def:g10:functions:variations}{decreasing} on $\intoo{-\infty}{0}$ and on
        $\intoo{0}{+\infty}$ (but \emph{not} on their \omterm{def:g10:numbers:interunion}{union}).
  \item $\abs x$ is strictly \omterm{def:g10:functions:variations}{decreasing} on $\intoc{-\infty}{0}$ and
        strictly \omterm{def:g11:func:monotone}{increasing} on $\intco{0}{+\infty}$.
\end{enumerate}
\end{proposition}

\begin{proof}
Let $u < v$. We factor $f(v) - f(u)$ in each case.

\emph{1.} $v^2 - u^2 = (v-u)(v+u)$. On $\intco{0}{+\infty}$: $v - u > 0$
and $v + u > 0$ (as $v > u \geq 0$), so $v^2 > u^2$. On
$\intoc{-\infty}{0}$: $v + u < 0$, so $v^2 < u^2$.

\emph{2.} $v^3 - u^3 = (v - u)(v^2 + uv + u^2)$, and the second factor
equals $\left(v + \frac u2\right)^2 + \frac{3u^2}{4} > 0$ unless
$u = v = 0$. Hence $v^3 > u^3$ whenever $v > u$.

\emph{3.} For $0 \leq u < v$:
$\sqrt v - \sqrt u = \frac{(\sqrt v - \sqrt u)(\sqrt v + \sqrt u)}
{\sqrt v + \sqrt u} = \frac{v - u}{\sqrt v + \sqrt u} > 0$.

\emph{4.} For $0 < u < v$:
$\frac1v - \frac1u = \frac{u - v}{uv} < 0$ since $u - v < 0$ and $uv > 0$;
same computation for $u < v < 0$. On the \omterm{def:g10:numbers:interunion}{union} the claim fails:
$-1 < 1$ but $\frac{1}{-1} = -1 < 1 = \frac11$.

\emph{5.} On $\intco{0}{+\infty}$, $\abs x = x$ is strictly \omterm{def:g11:func:monotone}{increasing}; on
$\intoc{-\infty}{0}$, $\abs x = -x$ is strictly \omterm{def:g10:functions:variations}{decreasing}.
\end{proof}

\begin{remark}
Point 4 is a classic trap: ``\omterm{def:g10:functions:variations}{decreasing} on $\intoo{-\infty}{0}$ and on
$\intoo{0}{+\infty}$'' is \emph{not} the same as ``\omterm{def:g10:functions:variations}{decreasing} on
$\R \setminus \{0\}$''. Variations only make sense on \omterm{def:g10:numbers:interval}{intervals}.
\end{remark}

\begin{example}[Comparing without computing]
Since $\sqrt x$ is strictly \omterm{def:g11:func:monotone}{increasing}, $\sqrt{7} < \sqrt{8}$; since
$x^2$ is strictly \omterm{def:g10:functions:variations}{decreasing} on $\intoc{-\infty}{0}$,
$(-3)^2 > (-2)^2$. Monotonicity turns comparisons of \omterm{def:g10:functions:function}{images} into
comparisons of arguments.
\end{example}

\section{New functions from old}

\begin{proposition}[Transformations]\label{prop:g11:func:transforms}
Let $u$ be a \omterm{def:g11:func:function}{function} defined on an \omterm{def:g10:numbers:interval}{interval} $I$ and $k \in \R$.
\begin{enumerate}
  \item $u + k$ has the \emph{same} variations as $u$ on $I$; its curve is
        that of $u$ shifted vertically by $k$.
  \item For $\lambda > 0$, $\lambda u$ has the same variations as $u$; for
        $\lambda < 0$, the variations are reversed.
  \item If $u \geq 0$ on $I$, then $\sqrt u$ has the same variations as
        $u$.
  \item If $u \neq 0$ and $u$ has constant sign on $I$, then $\frac1u$ has
        the variations \emph{opposite} to those of $u$.
\end{enumerate}
\end{proposition}

\begin{proof}
Let $u < v$ in $I$ (abusing notation, we compare $x$-values $u < v$; write
$s = u(x_1)$, $t = u(x_2)$ for the \omterm{def:g10:functions:function}{images} of $x_1 < x_2$).

\emph{1.} $(u+k)(x_2) - (u+k)(x_1) = t - s$: the difference of \omterm{def:g10:functions:function}{images} is
unchanged, so its sign is unchanged.

\emph{2.} $\lambda t - \lambda s = \lambda(t - s)$: the sign is kept if
$\lambda > 0$, flipped if $\lambda < 0$.

\emph{3.} If $u$ is \omterm{def:g11:func:monotone}{increasing} and $x_1 < x_2$, then $0 \leq s \leq t$, and
$\sqrt{}$ being \omterm{def:g11:func:monotone}{increasing} (\cref{prop:g11:func:refvariations}) gives
$\sqrt s \leq \sqrt t$. Same argument in the \omterm{def:g10:functions:variations}{decreasing} case.

\emph{4.} If $u$ is \omterm{def:g11:func:monotone}{increasing}, $0 < s \leq t$ (or $s \leq t < 0$), then
$\frac1s \geq \frac1t$ because $x \mapsto \frac1x$ is \omterm{def:g10:functions:variations}{decreasing} on each
side of $0$: the order is reversed.
\end{proof}

\begin{example}\label{ex:g11:func:composite}
Study $f(x) = \sqrt{x^2 + 1}$ on $\R$. The inner \omterm{def:g11:func:function}{function}
$u(x) = x^2 + 1$ is positive, \omterm{def:g10:functions:variations}{decreasing} on $\intoc{-\infty}{0}$ and
\omterm{def:g11:func:monotone}{increasing} on $\intco{0}{+\infty}$ (a shifted square). By point 3, $f$
has the same variations: \omterm{def:g10:functions:variations}{decreasing} then \omterm{def:g11:func:monotone}{increasing}, with \omterm{def:g10:functions:extrema}{minimum}
$f(0) = 1$.
\end{example}

\begin{omfigure}
\begin{tikzpicture}
\begin{axis}[omaxis,
  width=8.6cm, height=5.8cm,
  xmin=-2.6, xmax=2.6, ymin=-1.7, ymax=4.6,
  xtick={-1,1}, ytick={1,2}]
  \addplot[omDef, thick, domain=-1.9:1.9] {x^2};
  \addplot[omThm, thick, domain=-1.55:1.55] {x^2 + 1};
  \addplot[omProp, thick, domain=-2.4:2.4, samples=80]
    {sqrt(x^2 + 1)};
  \node[omDef, font=\small, anchor=west] at (axis cs:1.9,3.4) {$x^2$};
  \node[omThm, font=\small, anchor=west] at (axis cs:1.35,3.3)
    {$x^2{+}1$};
  \node[omProp, font=\small, anchor=north] at (axis cs:-1.9,1.95)
    {$\sqrt{x^2+1}$};
\end{axis}
\end{tikzpicture}

{\small Transformations in action: shifting $x^2$ up by $1$ (red) keeps
its variations; taking the square root (orange) keeps them again, as in
\cref{ex:g11:func:composite}.}
\end{omfigure}

\section{Even and odd functions}

\begin{definition}[Parity]\label{def:g11:func:parity}
Let $f$ be defined on a \omterm{def:g11:func:function}{domain} $D$ symmetric about $0$ (i.e.\ $-x \in D$
whenever $x \in D$). The \omterm{def:g11:func:function}{function} $f$ is \emph{even}\index{even function}
if $f(-x) = f(x)$ for all $x \in D$, and \emph{odd}\index{odd function} if
$f(-x) = -f(x)$ for all $x \in D$.
\end{definition}

\begin{proposition}[Geometric meaning]\label{prop:g11:func:paritysym}
The curve of an \omterm{def:g11:func:parity}{even function} is symmetric about the $y$-axis; the curve
of an \omterm{def:g11:func:parity}{odd function} is symmetric about the origin.
\end{proposition}

\begin{proof}
If $f$ is \omterm{def:g11:func:parity}{even} and $(x, y)$ is on the curve (so $y = f(x)$), then
$f(-x) = f(x) = y$: the mirror point $(-x, y)$ is also on the curve. If
$f$ is \omterm{def:g11:func:parity}{odd}, $f(-x) = -y$: the point $(-x, -y)$, symmetric of $(x,y)$ about
the origin, is on the curve.
\end{proof}

\begin{example}
$x^2$ and $\abs x$ are \omterm{def:g11:func:parity}{even}; $x^3$ and $\frac1x$ are \omterm{def:g11:func:parity}{odd};
$f(x) = x^2 + x$ is neither: $f(-1) = 0$ while $f(1) = 2$, so
$f(-1) \neq \pm f(1)$. Parity halves the work: an \omterm{def:g11:func:parity}{even} or \omterm{def:g11:func:parity}{odd function}
only needs to be studied on $\intco{0}{+\infty}$.
\end{example}

\section{Exercises}

\begin{exercise}[$\star$]\label{exo:g11:func:1}
Give the \omterm{def:g11:func:function}{domain} of each \omterm{def:g11:func:function}{function}:
\[
f(x) = \sqrt{3 - x}, \qquad
g(x) = \frac{1}{x^2 - 4}, \qquad
h(x) = \frac{\sqrt{x}}{x - 5} .
\]
\end{exercise}

\begin{exercise}[$\star$]\label{exo:g11:func:2}
Using the variations of the \omterm{def:g11:func:reference}{reference functions}, compare without a
calculator:
\[
\sqrt{11} \text{ and } \sqrt{13}; \qquad
(-2.1)^2 \text{ and } (-2.05)^2; \qquad
\frac{1}{0.9} \text{ and } \frac{1}{1.1}; \qquad
(-1.01)^3 \text{ and } (-0.99)^3 .
\]
\end{exercise}

\begin{exercise}[$\star$]\label{exo:g11:func:3}
Determine whether each \omterm{def:g11:func:function}{function} is \omterm{def:g11:func:parity}{even}, \omterm{def:g11:func:parity}{odd}, or neither:
\[
f(x) = 3x^4 - x^2, \qquad
g(x) = x^3 + x, \qquad
h(x) = x^2 + x^3, \qquad
k(x) = \frac{x}{x^2+1} .
\]
\end{exercise}

\begin{exercise}[$\star$]\label{exo:g11:func:4}
A \omterm{def:g11:func:function}{function} $f$ has the following variation behaviour: \omterm{def:g11:func:monotone}{increasing} on
$\intcc{-3}{1}$ from $f(-3) = -2$ to $f(1) = 4$, then \omterm{def:g10:functions:variations}{decreasing} on
$\intcc{1}{5}$ down to $f(5) = 0$. How many solutions does the \omterm{def:g10:algebra:equation}{equation}
$f(x) = 1$ have on $\intcc{-3}{5}$? And $f(x) = 5$?
\end{exercise}

\begin{exercise}[$\star\star$]\label{exo:g11:func:5}
Using \cref{met:g11:func:variation} (sign of $f(v) - f(u)$), prove that
$f(x) = x^2 - 4x$ is strictly \omterm{def:g11:func:monotone}{increasing} on $\intco{2}{+\infty}$.
\end{exercise}

\begin{exercise}[$\star\star$]\label{exo:g11:func:6}
Without computing any derivative, give the variations of
\[
f(x) = \sqrt{5 - x} \ \text{ on } \intoc{-\infty}{5},
\qquad
g(x) = \frac{1}{x^2 + 1} \ \text{ on } \intco{0}{+\infty},
\qquad
h(x) = 3 - 2\sqrt{x} \ \text{ on } \intco{0}{+\infty}.
\]
\end{exercise}

\begin{exercise}[$\star\star$]\label{exo:g11:func:7}
Let $f(x) = (x-2)^2 - 3$ on $\R$.
\begin{enumerate}
  \item Give the variations of $f$ and its \omterm{def:g10:functions:extrema}{minimum}.
  \item Deduce the variations of $g(x) = \frac{1}{(x-2)^2+1}$ and its
        \omterm{def:g10:functions:extrema}{maximum}. (Note $g = \frac{1}{f + 4}$.)
\end{enumerate}
\end{exercise}

\begin{exercise}[$\star\star$]\label{exo:g11:func:8}
The curve of a \omterm{def:g11:func:function}{function} $f$ passes through the points $(0, 1)$, $(2, 3)$
and $(4, 1)$, and $f$ is \omterm{def:g11:func:monotone}{increasing} on $\intcc{0}{2}$ and \omterm{def:g10:functions:variations}{decreasing} on
$\intcc{2}{4}$. Sketch a possible curve, then sketch the curves of
$f + 2$, of $-f$, and of $2f$.
\end{exercise}

\begin{exercise}[$\star\star$]\label{exo:g11:func:9}
Show that the \omterm{def:g11:func:function}{function} $f(x) = \frac{2x+1}{x+1}$ can be written
$f(x) = 2 - \frac{1}{x+1}$, and deduce its variations on
$\intoo{-1}{+\infty}$.
\end{exercise}

\begin{exercise}[$\star\star$]\label{exo:g11:func:10}
True or false (justify or give a counterexample):
\begin{enumerate}
  \item The sum of two \omterm{def:g11:func:monotone}{increasing functions} on $I$ is \omterm{def:g11:func:monotone}{increasing} on $I$.
  \item The product of two \omterm{def:g11:func:monotone}{increasing functions} on $I$ is \omterm{def:g11:func:monotone}{increasing} on
        $I$.
  \item If $f$ is \omterm{def:g11:func:monotone}{increasing} on $I$, then $f^2$ ($= f \times f$) is
        \omterm{def:g11:func:monotone}{increasing} on $I$.
\end{enumerate}
\end{exercise}

\begin{exercise}[$\star\star\star$]\label{exo:g11:func:11}
Let $f(x) = x + \frac{1}{x}$ for $x > 0$.
\begin{enumerate}
  \item Show that for $0 < u < v$,
        $f(v) - f(u) = (v - u)\,\frac{uv - 1}{uv}$.
  \item Deduce that $f$ is strictly \omterm{def:g10:functions:variations}{decreasing} on $\intoc{0}{1}$ and
        strictly \omterm{def:g11:func:monotone}{increasing} on $\intco{1}{+\infty}$.
  \item Conclude that $x + \frac1x \geq 2$ for all $x > 0$, with equality
        only at $x = 1$.
\end{enumerate}
\end{exercise}
